\documentclass[../../main.tex]{subfiles}
\begin{document}

{\bf[解]}

\subparagraph{(1)}

% TODO: 3通りの右端タイルの埋め方
\begin{figure}[htb]
  \centering
  \begin{tikzpicture}[scale=0.9, >=stealth, every node/.style={font=\small}]

    % ============================================================
    % パターン 1: 右端に 1×2 (縦置き) タイルを 1 個置く (幅 1)
    % 残り: 幅 n-1 の敷き詰め -> A_{n-1} 通り
    % ============================================================
    \begin{scope}[yshift=4.6cm]
      % 残りの領域 (幅 n-1): 薄いグレー＋ハッチング
      \fill[gray!12] (0, 0) rectangle (4.5, 2.0);
      \draw[dashed, gray!60] (0, 0) rectangle (4.5, 2.0);
      \node[font=\normalsize\bfseries, gray!75!black] at (2.25, 1.0) {$A_{n-1}$ 通り};

      % 右端のタイル: 1×2 (縦置き)
      \fill[blue!20] (4.5, 0) rectangle (5.5, 2.0);
      \draw[very thick, blue!85!black] (4.5, 0) rectangle (5.5, 2.0);
      \node[blue!90!black, font=\scriptsize, align=center] at (5.0, 1.0) {縦\\$1\times 2$};

      % 全体枠線
      \draw[thick] (0, 0) rectangle (5.5, 2.0);

      % 寸法ブレース
      \draw[decorate, decoration={brace, amplitude=4pt}, thick, gray!70]
      (0, 2.15) -- (4.5, 2.15) node[midway, above=4pt, font=\footnotesize] {$n-1$};
      \draw[decorate, decoration={brace, amplitude=3pt}, thick, blue!85!black]
      (4.5, 2.15) -- (5.5, 2.15) node[midway, above=4pt, font=\footnotesize, blue!90!black] {$1$};
      \draw[decorate, decoration={brace, mirror, amplitude=4pt}, thick, gray!70]
      (-0.15, 0) -- (-0.15, 2.0) node[midway, left=4pt, font=\footnotesize] {$2$};

      % パターン見出し
      \node[left=18pt, font=\footnotesize\bfseries, align=right] at (-0.15, 1.0)
      {ケース 1\\[1pt]\scriptsize (右端幅 $1$)};
    \end{scope}

    % ============================================================
    % パターン 2: 右端に 2×1 (横置き) タイルを 2 個置く (幅 2)
    % 残り: 幅 n-2 の敷き詰め -> A_{n-2} 通り
    % ============================================================
    \begin{scope}[yshift=2.3cm]
      % 残りの領域 (幅 n-2)
      \fill[gray!12] (0, 0) rectangle (3.5, 2.0);
      \draw[dashed, gray!60] (0, 0) rectangle (3.5, 2.0);
      \node[font=\normalsize\bfseries, gray!75!black] at (1.75, 1.0) {$A_{n-2}$ 通り};

      % 右端のタイル: 2×1 (横置き 2 段)
      \fill[teal!20] (3.5, 0) rectangle (5.5, 2.0);
      \draw[very thick, teal!85!black] (3.5, 0) rectangle (5.5, 1.0);
      \draw[very thick, teal!85!black] (3.5, 1.0) rectangle (5.5, 2.0);
      \node[teal!90!black, font=\scriptsize] at (4.5, 1.5) {横 $2\times 1$};
      \node[teal!90!black, font=\scriptsize] at (4.5, 0.5) {横 $2\times 1$};

      % 全体枠線
      \draw[thick] (0, 0) rectangle (5.5, 2.0);

      % 寸法ブレース
      \draw[decorate, decoration={brace, amplitude=4pt}, thick, gray!70]
      (0, 2.15) -- (3.5, 2.15) node[midway, above=4pt, font=\footnotesize] {$n-2$};
      \draw[decorate, decoration={brace, amplitude=4pt}, thick, teal!85!black]
      (3.5, 2.15) -- (5.5, 2.15) node[midway, above=4pt, font=\footnotesize, teal!90!black] {$2$};

      % パターン見出し
      \node[left=18pt, font=\footnotesize\bfseries, align=right] at (-0.15, 1.0)
      {ケース 2\\[1pt]\scriptsize (右端幅 $2$)};
    \end{scope}

    % ============================================================
    % パターン 3: 右端に 2×2 (正方形) タイルを 1 個置く (幅 2)
    % 残り: 幅 n-2 の敷き詰め -> A_{n-2} 通り
    % ============================================================
    \begin{scope}[yshift=0.0cm]
      % 残りの領域 (幅 n-2)
      \fill[gray!12] (0, 0) rectangle (3.5, 2.0);
      \draw[dashed, gray!60] (0, 0) rectangle (3.5, 2.0);
      \node[font=\normalsize\bfseries, gray!75!black] at (1.75, 1.0) {$A_{n-2}$ 通り};

      % 右端のタイル: 2×2 (正方形)
      \fill[orange!20] (3.5, 0) rectangle (5.5, 2.0);
      \draw[very thick, orange!90!black] (3.5, 0) rectangle (5.5, 2.0);
      \node[orange!95!black, font=\scriptsize, align=center] at (4.5, 1.0) {正方形\\$2\times 2$};

      % 全体枠線
      \draw[thick] (0, 0) rectangle (5.5, 2.0);

      % 寸法ブレース
      \draw[decorate, decoration={brace, amplitude=4pt}, thick, gray!70]
      (0, 2.15) -- (3.5, 2.15) node[midway, above=4pt, font=\footnotesize] {$n-2$};
      \draw[decorate, decoration={brace, amplitude=4pt}, thick, orange!90!black]
      (3.5, 2.15) -- (5.5, 2.15) node[midway, above=4pt, font=\footnotesize, orange!95!black] {$2$};

      % パターン見出し
      \node[left=18pt, font=\footnotesize\bfseries, align=right] at (-0.15, 1.0)
      {ケース 3\\[1pt]\scriptsize (右端幅 $2$)};
    \end{scope}

  \end{tikzpicture}
  \caption{右端のタイルの敷き詰め方による場合分け}
  \label{fig:tiling_recurrence_cases}
\end{figure}

$A_n$の時の右端のタイルの敷き詰め方で場合分けして考える．
右端1または2マスの敷き詰め方は図の3通りであり，それぞれ対応する並べ方は$A_{n-1},A_{n-2},A_{n-2}$だから
\begin{align}
  A_n=A_{n-1}+2A_{n-2}\cdots\text{(答)} \label{eq:1}
\end{align}
である．

\subparagraph{(2)}

\cref{eq:1}を変形すると
\begin{align}
  \begin{cases}
    A_{n+2}-2A_{n+1}=-(A_{n+1}-2A_{n}) \\
    A_{n+2}+A_{n+1}=2( A_{n+1}+A_{n})
  \end{cases}
\end{align}
これを繰り返し用いて，初期条件$A_1=1$，$A_2=3$から
\begin{align}
  \begin{cases}
    A_{n+1}-2A_{n}=(-1)^{n-1}(3-2) \\
    A_{n+1}+A_{n}=2^{n-1}(3+1)
  \end{cases}
\end{align}
辺々引いて整理すれば
\begin{align}
  A_n=\dfrac{1}{3}\left(2^{n+1}+(-1)^n\right)\cdots\text{(答)}
\end{align}
を得る．
\newpage
\end{document}
